\documentclass{standalone}
\usepackage{circuitikz}
\usetikzlibrary{calc}
\definecolor{circuitnotemark}{HTML}{FE00FE}
\begin{document}
\begin{circuitikz}[american, line width=0.8pt]
\ctikzset{bipoles/length=1.2cm}
\coordinate (b3) at (4,-2);
\coordinate (b7) at (12,-2);
\coordinate (d7) at (12,-6);
\coordinate (d9) at (16,-6);
\coordinate (h3) at (4,-14);
\coordinate (h5) at (8,-14);
\coordinate (j3) at (4,-18);
\coordinate (j7) at (12,-18);
\coordinate (f7) at (12,-10);
\coordinate (f9) at (16,-10);
\draw (b3) to[TL, l_=$\mathrm{L1}$, a^=$50\,\Omega$] (b7); % line 3
\draw (d7) to[C, l_=$C_{1}$, a^=$100\,\mathrm{p}\mathrm{F}$] (d9); % line 4
\draw (h3) to[C, l_=$C_{2}$] (h5); % line 5
\draw (j3) to[TL, l_=$\mathrm{L2}$, a^=$146\,\Omega$] (j7); % line 6
\draw (b7) -- (d7); % line 8
\draw (d7) -- (f7); % line 9
\draw (f7) -- (f9); % line 10
\draw (f9) -- (d9); % line 11
\node[circ] at (d7) {};
\node[anchor=south west, inner sep=2pt, circuitnotemark, font=\LARGE\bfseries] (circuittitle) at (current bounding box.north west) {X};
\path (circuittitle.south west) rectangle ++(12.619,0.853);
\end{circuitikz}
\end{document}
